% !TeX root = ../main.tex
% !Mode:: "TeX:UTF-8"

\chapter*{结\quad 论}

本文面向 6G 物联网中多源状态更新的实时采集与调度问题，以信息新鲜度（Age of Information，AoI）为核心性能
指标，围绕"如何在有限资源下高效采集数据并保障信息及时性"这一主题展开研究。论文将深度学习特征压缩与李雅普诺夫
优化调度相结合，提出了一套完整的系统建模与算法设计方法，并通过仿真实验验证了其有效性。主要工作与创新点
总结如下：

\begin{enumerate}
  \item \textbf{建立了多源状态更新与排队传输的系统模型。}针对多源传感节点经队列汇聚至基站的典型结构，推导了
        AoI 与队列长度的演化方程，并在队列稳定约束下，将数据采集与调度问题形式化为最小化长时平均加权 AoI 的
        随机优化问题，为后续方法设计提供了清晰的优化目标与约束框架。
  \item \textbf{提出了基于堆叠自编码器（SAE）的状态数据压缩方法。}针对高维传感数据直接传输开销过大的问题，
        利用 SAE 进行无监督逐层预训练，学习数据的低维特征表示。实验表明，在压缩率为 $4\times$ 时重建均方误差
        仅为 $0.058$，能够在较小精度损失下将传输开销降低至 $25\%$，为大规模多源采集提供了可行途径。
  \item \textbf{设计了基于李雅普诺夫漂移最小化的在线调度算法。}算法在每个时隙仅依赖当前可观测的队列状态与
        AoI 信息，通过最小化"漂移加惩罚"项选择最优源进行传输，无需未来信道或到达统计信息，每时隙复杂度仅为
        $O(N)$。理论分析表明，所提算法可在保证队列速率稳定的同时，使长时平均加权 AoI 以 $O(1/V)$ 的速率逼近
        最优值。
  \item \textbf{通过仿真验证了所提方法的性能优势。}与轮询、最大 AoI 优先、最大队列优先等基准算法相比，
        所提方法在多种系统负载下均取得最低的平均加权 AoI（较最大 AoI 优先降低约 $28\%$），且队列长度保持
        适中，充分体现了其在时效性与稳定性之间的良好平衡。
\end{enumerate}

展望未来的工作，以下几个方向值得进一步深入研究：其一，本文仅考虑了单服务队列系统，未来可将模型推广至多信道
并行传输、面向更复杂时变信道情形下的系统；其二，可将 SAE 压缩与调度决策进行端到端联合优化，在训练阶段即引入
时效性目标，进一步挖掘"压缩—传输—调度"耦合带来的增益；其三，如何将所提方法部署到真实物联网平台，并考虑能量
约束、设备异构等工程因素，也是具有重要意义的研究方向。

此外，随着深度学习与在线优化技术的不断发展，将强化学习等自适应方法引入多源调度，有望在无需显式建模系统动态的
前提下实现更优的时效性能，这也是后续研究值得探索的重要路径。
